\section{Introduction}\label{sec:introduction}
Vision-language-action (VLA) models are usually adapted to a new robot by supervised fine-tuning on demonstrations collected with that robot. SmolVLA, for example, fine-tunes on 50 demonstrations per real-world task~\cite{smolvla}. iRe-VLA begins its real-robot experiments with 2,000 teleoperated trajectories~\cite{irevla}, and PLD fine-tunes its base policy on 200 teleoperated trajectories before residual RL~\cite{pld}. Each new embodiment needs a teleoperation interface, operator time, and a fresh collection campaign. We investigate whether this collection step can be avoided by making reinforcement learning (RL) the first stage of target-embodiment adaptation.

RL from a pretrained VLA on a new action interface faces three obstacles. The initial policy rarely completes a multi-stage task, so binary success provides little guidance. Training component skills separately helps exploration, but later skills can erase earlier ones. Finally, the retained skills must be composed into one behavior for the complete instruction. Consider \emph{put the apple on the plate}. A failed rollout may contain a correct approach and grasp followed by a poor release, and treating the entire attempt as a failure discards that distinction.

We address these obstacles with the schedule in Fig.~\ref{fig:hero}. A newly initialized residual controller is placed over a retained pretrained SmolVLA action expert~\cite{smolvla}. It first acquires approach, pickup, and placement skills through RL with dense rewards. Binary refinement then targets skill completion, and instruction-balanced self-imitation retains successful behavior across instructions. Finally, milestone-based RL trains the same policy to execute the complete instruction from an empty-gripper start.

The data requirement is precise. We collect no human demonstrations on the target robot. Supervised stages fit only automatically generated trajectories: successful policy rollouts, a scripted simulator oracle, and policy rollouts with scripted controller assistance (Section~\ref{sec:lineage}). These trajectories are demonstrations, so the setting is not demonstration-free; it removes human collection. Simulator construction, reward and reset design, oracle engineering, and RL compute remain. We therefore make no claim of lower total adaptation cost.

This paper extends our earlier preprint, which adapted OpenVLA-OFT to the same cable-driven parallel robot (CDPR) with RL alone and reported directional motion and object reaching, with 9.75\% strict success on \emph{move to \textless object\textgreater}~\cite{nurtdinov2026rlboot}. Here we address complete pick-and-place, skill consolidation, and a different controller, a residual over SmolVLA. A separate study on LIBERO-Spatial with OpenVLA-OFT~\cite{openvlaoft,libero} examines intermediate feedback for full-task RL. It does not evaluate the full schedule.

Our contributions are:
\begin{enumerate}
\item A staged adaptation method for a pretrained VLA on a new robot without target-embodiment human demonstrations: dense-reward skill acquisition by a new residual controller, binary refinement with instruction-balanced self-imitation, and milestone-based learning of the complete instruction.
\item A milestone credit rule that retains local credit for achieved approach and pickup when later placement fails.
\item An empirical instantiation on a simulated CDPR, with a full training-lineage account, a skill-by-stage record, and strict full-task evaluation; and an auxiliary LIBERO-Spatial comparison of intermediate feedback.
\end{enumerate}
